\SourceLeaf{321}%
% A544-P0321-D01
\section*{BEMERKUNGEN.}
\pdfbookmark[1]{Bemerkungen zur geometrischen Seite der ternären Formen}{ternary-bemerkungen}
\BookRule{.15\textwidth}
Neben den vorstehenden Notizen, welche die in der Anzeige von \textsc{Seeber}'s Untersuchungen der ternären Formen gegebenen Gesichtspunkte theilweise weiter entwickeln, sind in der Handschrift mehre eigne mit einem achtzölligen \textsc{Reichenbach}'schen Theodolithen ausgeführte Crystallmessungen aufgezeichnet. Die einzelnen Protokolle enthalten das jedesmalige Datum der Beobachtung, woraus zu ersehen ist, dass diese Untersuchung dem Monat Juli $1831$ angehört.

% A544-P0321-D02
Aus der Theorie der indifferenten ternären quadratischen Formen findet sich im handschriftlichen Nachlass nur der folgende, wahrscheinlich in der Zeit der Ausarbeitung der Disqu. Arr. aufgezeichnete Lehrsatz
\begin{quote}
‘Omnes transformationes formae ternariae
\[
\left(\begin{smallmatrix}1,&1,&-1\\0,&0,&0\end{smallmatrix}\right)
\]
in se ipsam exhibentur per formulam
\[
\begin{array}{ccc}
\alpha\delta+\beta\gamma&\alpha\beta-\gamma\delta&\alpha\beta+\gamma\delta\\[.5ex]
\alpha\gamma-\beta\delta&\frac12(\alpha\alpha+\delta\delta-\beta\beta-\gamma\gamma)&\frac12(\alpha\alpha+\gamma\gamma-\beta\beta-\delta\delta)\\[.5ex]
\alpha\gamma+\beta\delta&\frac12(\alpha\alpha+\beta\beta-\gamma\gamma-\delta\delta)&\frac12(\alpha\alpha+\beta\beta+\gamma\gamma+\delta\delta)
\end{array}
\]
acceptis $\alpha,\beta,\gamma,\delta$ ita ut fiat $\alpha\delta-\beta\gamma=1$.’
\end{quote}

% A544-P0321-D03
Es entstehen nemlich alle Transformationen, in denen die neun Coëfficienten ganze Zahlen sind, wenn für $\alpha,\beta,\gamma,\delta$ sowohl alle die der Bedingungsgleichung genügenden ganzen Zahlen und zwar zwei gerade und zwei ungerade gesetzt werden, als auch alle die ungeraden Vielfache von $\surd\frac12$, welche dieselbe Bedingungsgleichung $\alpha\delta-\beta\gamma=1$ erfüllen.

\SourceLeaf{322}%
% A544-P0322-D01
Zu Seite $309$. Chaux carbonatée équiaxe, inverse, contrastante und mixte sind die von \textsc{Hauy} (Traité de Minéralogie $1801$ Tome II pag. $132$, $137$) gebrauchten Benennungen.

% A544-P0322-D02
Die Tafel der Transformationen der Form $\left(\begin{smallmatrix}5,&5,&5\\1,&1,&1\end{smallmatrix}\right)$ enthält in der ersten Verticalreihe die Coëfficienten der Substitution, in der zweiten die dadurch entstandene neue Form, in der dritten die der letztern Form entsprechende primitive, wenn diese nicht selbst schon eine solche ist, und in der vierten deren Adjuncta.
\begin{flushright}\textsc{Schering.}\end{flushright}
\BookRule{.30\textwidth}

\SourceLeaf{323}%
% A544-P0323-D01
\chapter*{ZUR THEORIE DER BIQUADRATISCHEN RESTE.}
\pdfbookmark[0]{Zur Theorie der biquadratischen Reste}{nachlass-biquadratic}
\renewcommand{\BookRunningTitle}{Zur Theorie der biquadratischen Reste}
\BookRule{.18\textwidth}
\begin{center}\textbf{[I.]}\end{center}
\pdfbookmark[1]{I}{nachlass-biquadratic-I}
\renewcommand{\Article}[1]{\par\medskip\Needspace{4\baselineskip}\pdfbookmark[2]{Art. #1}{nachlass-biquadratic-I-art-#1}\centerline{\textbf{#1.}}\nobreak\smallskip}
\Article{1}
Wir erweitern das Gebiet der höhern Arithmetik, indem wir darin auch die imaginären Grössen aufnehmen. Bei der gegenwärtigen Untersuchung nennen wir eine ganze imaginäre Zahl jede Grösse $x+iy$, wenn $x,y$ reelle ganze Zahlen sind.

% A544-P0323-D02
\Article{2}
Die unendliche Anzahl imaginärer ganzer Zahlen lässt sich am bequemsten durch Punkte in einer unbegrenzten Ebene sinnlich darstellen; wir nennen schlechthin denjenigen Punkt, dessen Abscisse $x$, die Ordinate $y$ ist, den Punkt $x+iy$, alle Punkte, die ganze Zahlen vorstellen, sollen Ganzepunkte heissen.

% A544-P0323-D03
\Article{3}
Um etwas bestimmtes festzusetzen, sollen die Abscissen immer auf der linken Seite positiv, die Ordinaten oben positiv sein.

% A544-P0323-D04
\Article{4}
Die gerade Linie von dem Punkte $x+iy$ zu dem Punkte $x'+iy'$ gezogen soll schlechtweg die gerade Linie $(x+iy,x'+iy')$ heissen, wir nehmen dabei zugleich, insofern es darauf ankommt, auf die Richtung Rücksicht und unterscheiden also die gerade Linie $x+iy,x'+iy'$ von der $x'+iy',x+iy$.

\SourceLeaf{324}%
% A544-P0324-D01
\Article{5}
Der Kürze wegen wollen wir imaginäre Grössen wie $x+iy$ auch durch einen einzigen Buchstaben bezeichnen, wie $z$.

% A544-P0324-D02
\Article{6}
Die Figur, welche durch die geraden Linien $zz',z'z'',z''z'''\ldots z^{n-1}z^n,z^nz$ begrenzt wird, nennen wir schlechtweg die Figur $zz'z''z'''\dotstwo z^n$. Wir schliessen dabei den Fall nicht aus, wo etwa einige dieser Linien einander schneiden.

% A544-P0324-D03
\Article{7}
Durch $S(z,z',z''\ldots z^n)$ bezeichnen wir allgemein die Summe von so vielen reellen ganzen Zahlen, als Ganzepunkte innerhalb der Figur liegen, indem wir für jeden Punkt, um den die Grenzlinie der Figur einmal, zweimal, dreimal u.s.w. herumgeht, die Zahl $\pm1$, $\pm2$, $\pm3$ etc. setzen; die obern Zeichen gelten, wenn die Grenzlinie den Punkt so umgibt, dass dieser auf der rechten Seite der Figur liegt, die untern im entgegengesetzten Fall. Schneiden sich also keine Seiten der Figur, so ist $S(z,z',z''\dotstwo)$ schlechthin die Anzahl der Punkte innerhalb der Figur, positiv oder negativ genommen.

% A544-P0324-D04
\Article{8}
Offenbar ist immer
\[
\begin{aligned}
S(z,z',z''\ldots z^n)&=S(z',z'',z'''\dotstwo z^n,z)\\
&=S(z'',z'''\dotstwo z')\text{ etc.}\\
&=-S(z^n,z^{n-1}\dotstwo z'',z',z)\\
&=-S(z^{n-1}\pdot z^{n-2}\dotstwo z',z,z^n)\text{ etc.}
\end{aligned}
\]

% A544-P0324-D05
\Article{9}
Wie es hiebei mit den auf der Grenzlinie selbst liegenden Punkten gehalten werden soll, muss noch näher bestimmt werden. Es gibt viele Fälle, wo auf der Grenzlinie gar keine ganze Punkte liegen können: dann ist keine Bestimmung nöthig. Liegen aber auf der Grenzlinie $zz'$ solche Punkte, so zeigen wir durch ein zwischen $z$ und $z'$ eingeschobenes $+$ an, dass diese Punkte so betrachtet werden sollen, als lägen sie rechts von der Grenzlinie, so wie durch ein $-$, als lägen sie links. Auch werden wir wol ein $0$ oder $\frac12$ einschieben, wodurch angedeutet werden soll, dass sie gar nicht oder nur mit dem halben Werthe auf
\SourceLeaf{325}%
% A544-P0325-D01
jeder Seite in Betracht gezogen werden sollen. Falls einer oder der andere der Punkte $z,z',z''$ etc. selbst ein Ganzepunkt, so wird er, wo nicht ausdrücklich das Gegentheil gesagt wird, gar nicht mitgezählt, als insofern er zugleich etwa als Nicht-Eckpunkt auch in Betracht kommt.

% A544-P0325-D02
\Article{10}
Lehrsätze. Wenn alle $z,z',z''$ etc. um eine und dieselbe Ganzezahl vermehrt werden, so bleibt das $S$ ungeändert.

% A544-P0325-D03
Wenn $i$ in $-i$ und jedes Bindezeichen ins entgegengesetzte verwandelt wird, so ändert $S$ bloss das Zeichen.
\[
\begin{aligned}
S(z,z',z''\ldots z^n)={}&S(z,u,u'\dotstwo u^n,z',z'',z^n)\\
&-S(z,u,u'\dotstwo u^n,z')\\
={}&S(z,u,u',u''\dotstwo u^n,z^m,z^{m+1}\dotstwo z^n)\\
&-S(z,u,u',u''\dotstwo u^n,z^m,z^{m-1}\dotstwo z',z)
\end{aligned}
\]
wo die Bindezeichen correspondiren müssen, aber zwischen den rückwärts laufenden Gliedern entgegengesetzt werden.

% A544-P0325-D04
Ist $\zeta$ eine ganze Zahl $=a+bi$, so ist, wenn die gegenüberliegenden Bindezeichen entgegengesetzt,
\[
S(z,z',z'+\zeta,z+\zeta)=[bx'-ay']-[bx-ay]
\]
Hiebei ist zu bemerken, dass wenn $bx'-ay'$ selbst eine ganze Zahl ist, diese für $[bx'-ay']$ angenommen werde, wenn das Bindezeichen zwischen $z'$ und $z'+\zeta$ $+$ ist, hingegen $1$ oder $\frac12$ weniger, wenn dieses Bindezeichen $-$ oder $\frac12$ ist; bei $bx-ay$ gilt das Umgekehrte.

% A544-P0325-D05
Uebrigens gilt die Formel nur für den Fall, wo $a$ und $b$ keinen gemeinschaftlichen Divisor haben; ist ihr grösster gemeinschaftlicher Theiler $=h$, so hat man dafür zu nehmen
\[
h\left[\frac{bx'-ay'}h\right]-h\left[\frac{bx-ay}h\right]
\]

% A544-P0325-D06
\Article{11}
Wenden wir uns nun näher zu unserm Gegenstande selbst. Wenn für den Modulus $m=a+bi$ die Zahlen $f,f',f''$ etc. so beschaffen sind, dass sie erstlich alle nach dem Modulus $m$ unter sich incongruent sind, zweitens aber jede ganze Zahl einer von ihnen nothwendig congruent sein muss, so nennen wir den
